% Generated draft: requires mathematical and source review.
\documentclass[11pt]{article}
\usepackage[margin=25mm]{geometry}
\usepackage{fontspec}
\setmainfont{Latin Modern Roman}
\usepackage{amsmath,amssymb,mathtools,unicode-math}
\setmathfont{Latin Modern Math}
\usepackage{enumitem,needspace,xurl,hyperref}
\setlength{\parindent}{0pt}
\setlength{\parskip}{5pt}
\setlength{\emergencystretch}{3em}
\newcommand{\N}{\mathbb{N}}
\newcommand{\Z}{\mathbb{Z}}
\newcommand{\Q}{\mathbb{Q}}
\newcommand{\R}{\mathbb{R}}
\newcommand{\C}{\mathbb{C}}
\newcommand{\F}{\mathbb{F}}
\newcommand{\A}{\mathbb{A}}
\newcommand{\PP}{\mathbb{P}}
\newcommand{\E}{\mathbb{E}}
\newcommand{\eps}{\varepsilon}
\newcommand{\dd}{\,\mathrm{d}}
\begin{document}
\section*{UnsolvedMath ID: 20001918}
\textbf{Source problem code:} catalogue entry
\subsection{Definitions and mathematical statement}
Babenko's Question \(8.7.1\) asks for a continuous or smooth family \((g_t)_{0\leq t\leq1}\) of sphere metrics allowing conical singularities, with \(g_0\) round, \(g_1\) the preceding example, and \(\sigma(g_t)\) decreasing. The aid reconstructs the missing context as the Calabi--Croke metric, the double of an equilateral Euclidean triangle, and the scale-invariant systolic area
\[
 \sigma(g)=\frac{\operatorname{Area}(S^2,g)}{\operatorname{sys}(g)^2},
\]
where \(\operatorname{sys}(g)\) is the shortest nonconstant closed-geodesic length under the chosen conical-geodesic convention. This is a geodesic systole on \(S^2\), not the length of a homotopically nontrivial loop. The precise topology on conical metrics and the word decreasing, strict versus nonincreasing, are not defined in the extracted sentence.

For the round metric \(\sigma=1/\pi\). For the reconstructed Calabi--Croke endpoint, with three cone angles \(2\pi/3\), the aid uses \(\sigma=1/(2\sqrt3)\). After constant-area normalization \(\operatorname{Area}=4\pi\), the question becomes a path with nondecreasing systole from
\[
 \operatorname{sys}(g_0)=2\pi\quad\text{to}\quad
 \operatorname{sys}(g_1)=\sqrt{8\pi\sqrt3}.
\]
Strict decrease of \(\sigma\) would correspond to strict increase of the normalized systole.

The reported smooth-end obstruction assumes that \(g_t\) is an ordinary smooth metric for \(t\) in some initial interval and is \(C^2\)-continuous there. If \(\sigma(g_t)\) is nonincreasing from the round endpoint, then on a smaller initial interval
\[
 \sigma(g_t)=1/\pi,
\]
and all these metrics are Zoll: their geodesics are closed with a common period in the stated equality case. Thus a strictly decreasing path starting immediately at the round metric is impossible in this smooth-through-the-end category. This local plateau does not rule out a nonincreasing conical path which later leaves that class, and no uniform plateau duration is specified.

The aid also reports that no continuous nonincreasing path can reach the Calabi--Croke endpoint if all interior metrics are smooth spheres of revolution, using the corresponding rotational systolic obstruction. This conclusion requires a path topology for which the relevant area and systole comparison passes to the endpoint; that topology and the exact rotational theorem are not provided. It is a restriction on the proposed ansatz, not a resolution for arbitrary conical metrics.

The canonical record ends ``A simpler version of this question is,'' without the simpler question. Its missing statement is not invented. The original general conical-path question is reported open despite these partial obstructions.
\subsection{Short English statement}
The reconstructed question seeks a nonincreasing systolic-area path from the round sphere to the doubled equilateral triangle. The aid forces a Zoll plateau near the round endpoint if the path initially remains smooth and \(C^2\)-continuous, and reports an obstruction to keeping all intermediate metrics smooth and rotationally symmetric. General conical paths remain open.
\subsection{Sources}
\begin{itemize}
\item Local UnsolvedMath registry, permanent ID \(20001918\), source code \texttt{catalogue entry}; canonical collection \texttt{aim-geometry-notes.json}, notes[255]. Source item \(8.7\), actual question \(8.7.1\), attributed to Babenko.
\item AIM workshop ``Geodesics,'' Burns--Matveev report ``Open problems and questions about geodesics'': \url{https://aimath.org/pastworkshops/geodesicsproblems.pdf}. The supplied verification reports inspecting Section \(8.7\), pages 18--19, and adjacent records to reconstruct the endpoint and invariant. The exact extraction is truncated before a simpler version. This session did not open the PDF.
\item Supplied research attempt \(1\), \texttt{partial\_result}, problem status at run ``open,'' rigor label ``proved,'' novelty confidence ``low.'' Registry status \texttt{partially\_solved}, classification \texttt{PARTIAL-PROGRESS}, updated 14 August 2026. Original remarks and literature are empty; no full authoritative references for the Zoll/local or rotational obstruction, proposer-field value, or proposed year are supplied.
\end{itemize}
\end{document}
